\section*{第 \ref{ch:g8:hormonal-communication} 章——激素调节}

\begin{solution}{exo:g8:hormonal-communication:1}
产物释放入血，而非经导管
排出。脑底调控\omterm{def:g8:hormonal-communication:glands}{腺体}：统领其他\omterm{def:g8:hormonal-communication:glands}{腺体}，启动
\omterm{def:g8:puberty:puberty}{青春期}。甲状腺：机体的节律设定者。
肾上腺：警觉\omterm{def:g8:hormonal-communication:glands}{激素}。胰腺：
管理血糖。（\omterm{def:g8:reproductive-systems:female}{卵巢}与\omterm{def:g8:reproductive-systems:male}{睾丸}：
调控周期与组织更新的\omterm{def:g8:hormonal-communication:glands}{激素}。）
\end{solution}

\begin{solution}{exo:g8:hormonal-communication:2}
运输呈广播式，寻址发生在接收端：只有
装配了该\omterm{def:g8:hormonal-communication:glands}{激素}受体的\omterm{def:g5:first-look-at-cells:cell}{细胞}才会响应。每户
都会收到这封信，只有配对的钥匙能将其打开。
\end{solution}

\begin{solution}{exo:g8:hormonal-communication:3}
心跳加快（见\cref{ch:g7:heart-and-circulation}）；
气道扩张（见\cref{ch:g7:lungs-and-gas-exchange}）；
\omterm{prop:g7:digestion-and-nutrients:absorption}{肝脏}储存的\omterm{prop:g7:organs-at-work:bill}{葡萄糖}释放入血，
\omterm{prop:g4:journey-of-food:absorb}{血液}改道流向（见\cref{ch:g7:digestion-and-nutrients}）；
\omterm{def:g3:bones-and-muscles:muscle}{肌肉}（见\cref{ch:g7:organs-at-work}）——\omterm{def:g1:the-five-senses:sense}{皮肤}苍白、
发麻，凑齐了这一整套反应。
\end{solution}

\begin{solution}{exo:g8:hormonal-communication:4}
分泌\omterm{def:g8:hormonal-communication:glands}{激素}的\omterm{def:g8:hormonal-communication:glands}{腺体}监测其信号的结果，并在
目标达成时减速——是热源式，而非水龙头式。示例：
血糖。血糖上升时由胰腺的一对\omterm{def:g8:hormonal-communication:glands}{激素}储存，
下降时则释放。
\end{solution}

\begin{solution}{exo:g8:hormonal-communication:5}
储存\omterm{def:g8:hormonal-communication:glands}{激素}的信号——未被发出，或未被
受体接收。此时\omterm{prop:g7:organs-at-work:bill}{葡萄糖}在进食后堆积，
超出\omterm{def:g7:removing-wastes:kidneys}{肾脏}的重吸收能力，并出现在
\omterm{def:g7:removing-wastes:kidneys}{尿液}中。
\end{solution}

\begin{solution}{exo:g8:hormonal-communication:6}
位于脑底部——调控\omterm{def:g8:hormonal-communication:glands}{腺体}，\omterm{def:g5:brain-in-command:system}{神经系统}
在上方，\omterm{def:g8:hormonal-communication:glands}{激素}在下方。例如：惊吓——\omterm{def:g8:nervous-communication:neuron}{神经}
信号先行，肾上腺分泌的\omterm{def:g8:hormonal-communication:glands}{激素}维持警觉；
或日照长度变化引发绵羊的繁殖\omterm{def:g8:hormonal-communication:glands}{激素}分泌。
\end{solution}

\begin{solution}{exo:g8:hormonal-communication:7}
水龙头要靠人拧开才会出水；热源会测量
自身的变化并自行调节。\omterm{def:g8:hormonal-communication:glands}{腺体}读取
其所调控的水平——分泌随需要升高或降低，
无需外部干预。
\end{solution}

\begin{solution}{exo:g8:hormonal-communication:8}
分泌者：胰腺，受血糖升高刺激。
接收者：首先是\omterm{prop:g7:digestion-and-nutrients:absorption}{肝脏}和\omterm{def:g3:bones-and-muscles:muscle}{肌肉}——它们在接收时储存
多余的\omterm{prop:g7:organs-at-work:bill}{葡萄糖}。\omterm{prop:g8:hormonal-communication:feedback}{反馈}：血糖下降本身
使分泌减弱。失效与应用：\omterm{ex:g8:hormonal-communication:diabetes}{糖尿病}——
以及药物的定量剂量，即借注射器送出的信号。
\end{solution}

\begin{solution}{exo:g8:hormonal-communication:9}
两者都是人工“写”出体内信号：注射器补充
进餐时补上胰腺缺失的储存信号；
\omterm{def:g8:contraception:pill}{避孕药}每天补充周期的“暂停”信号。
二者同属一种行为——模拟\omterm{def:g8:hormonal-communication:glands}{激素}，以有意
操控其系统——一个修复失效的通信，
另一个则主动覆盖健康的通信。
\end{solution}

\begin{solution}{exo:g8:hormonal-communication:10}
颤抖的起始：\omterm{def:g8:nervous-communication:neuron}{神经}性——瞬间的\omterm{def:g3:bones-and-muscles:muscle}{肌肉}
指令。胡须：\omterm{def:g8:hormonal-communication:glands}{激素}性——数月，遍及全身，
经\omterm{prop:g4:journey-of-food:absorb}{血液}传递。血糖通道：\omterm{def:g8:hormonal-communication:glands}{激素}性——\omterm{prop:g8:hormonal-communication:feedback}{反馈}
通信。膝跳反射：\omterm{def:g8:nervous-communication:neuron}{神经}性——一个\omterm{def:g8:nervous-communication:synapse}{突触}，
仅需五十分之一秒。
\end{solution}

\begin{solution}{exo:g8:hormonal-communication:11}
眼睛感知日照变短属于\omterm{def:g8:nervous-communication:neuron}{神经}调节；
脑底部的“办公室”将其转译，并统领
调控\omterm{def:g8:hormonal-communication:glands}{腺体}；其发出的信号唤醒\omterm{def:g8:reproductive-systems:female}{卵巢}自身的
\omterm{def:g8:hormonal-communication:glands}{激素}程序；母羊的周期开始——光线经由
“电报”传入，羔羊则经“邮政”送出。
\end{solution}

\begin{solution}{exo:g8:hormonal-communication:12}
论证：同一\omterm{prop:g4:journey-of-food:absorb}{血液}携带所有信号，但每种\omterm{def:g8:hormonal-communication:glands}{激素}只
在受体应答处起作用——听众决定
信道，因此广播式通信变为专线。推论：一个
\omterm{def:g5:first-look-at-cells:cell}{细胞}通过构建受体加入听众群体——因此
哪些组织接收信号本身也是逐器官构建的，
在发育过程中形成；\omterm{def:g8:puberty:puberty}{青春期}的\omterm{prop:g8:hormonal-communication:address}{靶器官}——喉、
\omterm{def:g1:the-five-senses:sense}{皮肤}、生长板——事先装配好了受体，
等待数年后接收第一封信号。
\end{solution}

\begin{solution}{pb:g8:hormonal-communication:1}
\textbf{1.} 胰腺成对的信号——储存
\omterm{def:g8:hormonal-communication:glands}{激素}在上升，其配对的释放\omterm{def:g8:hormonal-communication:glands}{激素}在下降，
——通过\omterm{prop:g8:hormonal-communication:feedback}{反馈}维持通道，机体则在\omterm{prop:g1:healthy-habits:needs}{睡眠}中。

\textbf{2.} 触发因素：进餐后的\omterm{prop:g7:organs-at-work:bill}{葡萄糖}经过
\omterm{prop:g7:digestion-and-nutrients:absorption}{绒毛}进入并提高血糖水平。分泌者：胰腺。
信号：储存\omterm{def:g8:hormonal-communication:glands}{激素}。接收者：\omterm{prop:g7:digestion-and-nutrients:absorption}{肝脏}和
\omterm{def:g3:bones-and-muscles:muscle}{肌肉}，它们储存多余的\omterm{prop:g7:organs-at-work:bill}{葡萄糖}。

\textbf{3.} 水龙头会越过目标，冲入
骤降之中；下降趋于平缓并进入通道，表明
分泌随其自身结果的到达而减弱——
能自我测量、自我停止。

\textbf{4.} 分泌量随水平升高而升高，随之回落：
分泌曲线紧随\omterm{prop:g7:organs-at-work:bill}{葡萄糖}曲线起伏，
峰后接峰，正是热源的走势。

\textbf{5.} 释放的一方（借助剧烈用力时
警报\omterm{def:g8:hormonal-communication:glands}{激素}的帮助）：\omterm{prop:g7:digestion-and-nutrients:absorption}{肝脏}所储存的
\omterm{prop:g7:organs-at-work:bill}{葡萄糖}——这笔储备在早餐时蓄满，
依\cref{ex:g7:digestion-and-nutrients:breadtrace}的
方式——被工作\omterm{def:g3:bones-and-muscles:muscle}{肌肉}的开销耗尽。

\textbf{6.} 同一释放的一方：
在无食的时段里稳定地打开储备，
与身体的消耗相匹配——
通道从另一侧维持。

\textbf{7.} 肾上腺。血糖为
\emph{战或逃}备好燃料：警报信件的全部要点，
在于让身体为即时用力而备足并改道，
无论用力是否真的发生。

\textbf{8.} 因为健康的信件往来对此作出了应答：
被释放的剩余\omterm{prop:g7:organs-at-work:bill}{葡萄糖}未被消耗（险情已过），
依次触发了储存信件，
通道随之\omterm{def:g7:organs-at-work:recovery}{恢复}——警报与热源先后相继。

\textbf{9.} 未经治疗时：早餐后的升高持续攀升，
不见回落——储存信件既未写出，也未听到——
上午维持高位平台；越过回收
阈值后，糖出现在\omterm{def:g7:removing-wastes:kidneys}{尿液}中
（\cref{ex:g7:removing-wastes:thrift}中的信号）。

\textbf{10.} 胰腺自己的餐时信件。
剂量须与餐食相匹配，因为信件的大小须
与待储存的\omterm{prop:g7:organs-at-work:bill}{葡萄糖}相称——过小则留下
堆积，过大则使通道向下越界，
这本身就是一场险情。

\textbf{11.} 夜间与下午：胰腺的一对，
\omterm{prop:g7:digestion-and-nutrients:absorption}{肝脏}与\omterm{def:g3:bones-and-muscles:muscle}{肌肉}予以应答——通道得以维持。晨餐与
晚餐：\omterm{prop:g7:organs-at-work:bill}{葡萄糖}发出信号，胰腺写出信件，
\omterm{prop:g7:digestion-and-nutrients:absorption}{肝脏}储存。中午：用力消耗，释放的一方
打开储备。下午五点：肾上腺插话，胰腺
随后收尾。同一个数值，五方发声，
全天没有一刻沉默。

\textbf{12.} “\omterm{prop:g8:hormonal-communication:feedback}{反馈}系统：读取自身信件结果的书写者，
这正是一个数值得以终生
沿狭窄通道行进、无人驾驭的
方式。”
\end{solution}
